\section{All-order quantization in the holomorphic local sector}
\label{sec:all-orders}

The order of a counterterm is its exponent of $\hbar$; it need not equal
the geometric loop number of a graph containing that counterterm.
Throughout this section the coefficients are complex, and the two sign
conventions $\epsilon,U\in\{1,-1\}$ are retained.  No choice of either
sign is needed for the existence theorem.

\subsection{The local contact stays in the required space}

The space $\Vht$ consists of integrated exterior products of holomorphic
jets of the superfield, with constant coefficients.  Integration uses the
full degree-three real and antiholomorphic exterior component.  In
particular, a ghost-number-zero element is cubic and a
ghost-number-one element is quartic.  Membership in $\Vht$ is a condition
on a local cochain, stronger than belonging to its image in cohomology.

\begin{lemma}[Closure of holomorphic local vertices under contacts]
\label{lem:full-wedge-contacts}
For the flat heat kernel on $\mathbb R\times\mathbb C^2$, every
total-collapse contact constructed from vertices in $\Vht$ belongs to
$\Vht$.  This includes vertices with arbitrary finite holomorphic
derivatives.  A contact of odd geometric loop number vanishes.
\end{lemma}

\begin{proof}
An integration by parts expresses every derivative at a vertex as a
finite sum of holomorphic derivatives on incident kernels and external
test fields.  These are the decorated graphs of
Wang--Williams~\cite[Propositions 3.18, 3.21--3.23]{WW}.  Their
compactified graph distributions are continuous on compactly supported
smooth test forms.  A nonzero total-collapse contact is supported on
the diagonal and differentiates its test form only holomorphically.

It remains to check its exterior structure.  Write $|V|$ and $|E|$ for
the numbers of vertices and internal edges of a connected collapsing
graph, and $n=|E|-|V|+1$ for its geometric loop number.  The tight
condition in three real and antiholomorphic exterior directions is
\begin{equation}
  2|E|=3|V|-4,
  \qquad |E|=3n+1,
  \qquad |V|=2n+2.
  \label{eq:contact-tight-degree}
\end{equation}
The extended kernel has exterior degree three on each edge.
After extracting the $|E|-1$ Schwinger differentials, its spatial
exterior degree is
\begin{equation}
  3|E|-(|E|-1)=2|E|+1=3(|V|-1).
  \label{eq:relative-form-saturation}
\end{equation}
All these differentials are differences between vertex coordinates.
They therefore saturate the relative real and antiholomorphic
directions.  Only the diagonal pullback of the external exterior
product remains.  No contraction of an external form index and no
additional background form is required.  The resulting local vertex
is precisely a full exterior product of holomorphic jets.  For cubic
inputs its external arity is
\begin{equation}
  3|V|-2|E|=4.
\end{equation}
Self-loops give zero: their differential factors are coordinate
differences, whose diagonal pullback remains zero after holomorphic
differentiation.  Finally, reflection of the real relative coordinates
changes the contact by $(-1)^n$.  Its local holomorphic test data are
unchanged.  This proves the odd-loop assertion.
\end{proof}

The regulated characters also record covariance and the limiting
step.  Let $r>0$ and
$|u_1|=|u_2|=1$, and pull back by
$(\tau,z_1,z_2)\mapsto(r\tau,ru_1z_1,ru_2z_2)$.  Direct substitution
in the normalized kernel gives
\begin{align}
  \phi_{r,u}^*[K_{r^2t}]_\Omega
    &=r^{-2}(u_1u_2)^{-1}[K_t]_\Omega,\notag\\
  \phi_{r,u}^*[D^*K_{r^2t}]_\Omega
    &=r^{-4}(u_1u_2)^{-1}[D^*K_t]_\Omega,\notag\\
  \phi_{r,u}^*P_{r^2\varepsilon}^{r^2L}
    &=r^{-2}(u_1u_2)^{-1}P_\varepsilon^L.
  \label{eq:regulated-covariance-all-orders}
\end{align}
These identities hold at positive cutoffs.  Continuity of pullback
on test forms passes them through the distributional boundary
restriction supplied by the compactified graph theorem.  Thus the
local obstruction has the prescribed rotation and dilation
characters.  Together with the exterior-form argument, this gives
the required statement at cochain level; exactness
in a larger space of local functionals alone would not prove it.

\subsection{The degree-four cohomology at every momentum degree}

We first work in the formal invariant trace algebra for traceless
matrices: distinct products of cyclic trace words are independent,
and a one-letter trace is zero.  For cohomology in degree four this
algebra agrees with the actual invariant tensors of $\fsl_N$ for
$N\ge5$, since only arities three, four, and five occur.  The
differential inserts a commutator and preserves the number of traces.

\begin{proposition}[The complete obstruction module]
\label{prop:all-order-h4}
For $w\ge2$, the formal trace complex satisfies
\begin{equation}
  H^3(C^0_{\fsl_N})_{\deg=2w}=0,
  \qquad
  H^4(C^0_{\fsl_N})_{\deg=2w}
       \cong\operatorname{Sym}^{2w-4}(\mathbb C^2).
  \label{eq:all-order-cohomology}
\end{equation}
The second space is represented entirely by two-trace cochains.
For the trace word $\tr(X_0X_1)\tr(X_2X_3)$, put
$P=p_0+p_1=-p_2-p_3$.  Its coefficient is
\begin{equation}
  [p_0P][p_2P]h(P),
  \qquad h\text{ homogeneous of degree }2w-4.
  \label{eq:all-order-double-trace}
\end{equation}
The other pairings are fixed by alternation.  Consequently
\begin{equation}
 \dim H^4(C^0_{\fsl_N})_{\deg=2w}=2w-3,
 \qquad
 \dim\bigl(H^4(C^0_{\fsl_N})_{\deg=2w}\bigr)^{SL_2}
 =\begin{cases}1,&w=2,\\0,&w\ge3.\end{cases}
 \label{eq:all-order-invariant-vanishing}
\end{equation}
\end{proposition}

\begin{proof}
\emph{One trace.}
A trace word of arity $k$ is a closed polygon with $k$ vertices:
write its momenta as $p_i=x_{i+1}-x_i$, with $x_k=x_0$.
The word coefficient is a polynomial in $x_0,\ldots,x_{k-1}$,
invariant under their simultaneous translation and cyclic with sign
$(-1)^{k-1}$.  Splitting one input into a commutator in the
Chevalley--Eilenberg differential adds two adjacent sides.  Equivalently,
the differential deletes the intervening polygon vertex.  Up to a
uniform degree sign this is the alternating deletion differential.

Here is the grading and acyclicity argument explicitly.  Let $j=k-1$.
Before imposing cyclicity, translation-invariant polynomials on
$j+1$ vertices have the ordinary additive bar differential on their
$j$ successive independent increments.  Its end faces omit the first
or last increment, and its interior faces add two consecutive
increments.  At fixed polynomial degree $m$, normalization removes
degenerate cochains.  The resulting complex is the graded dual of the
reduced bar complex of a polynomial algebra on two generators:
the coproduct $f(p+q)$ becomes multiplication after dualizing in the
factorial basis.  Write $HH$ for the cohomology of this full deletion
complex and $HC$ for the cohomology of its signed cyclic subcomplex.
The length-two Koszul resolution consequently gives
\begin{equation}
 HH^j_m=
 \begin{cases}
  \Lambda^j(\mathbb C^2),&m=j,\quad 0\le j\le2,\\
  0,&\text{otherwise}.
 \end{cases}
 \label{eq:ordinary-bar-bidegrees}
\end{equation}
This is a statement about translation-invariant polygon polynomials,
with no remaining total momentum variable.

The cyclic bicomplex is available degree by degree.  Rotation, each
face, and insertion of a zero side preserve $m$.  The complex with
its last face omitted is contracted by appending a duplicate final
vertex, with the alternating degree sign.  The finite cyclic norm
identifies cyclic invariants and coinvariants over characteristic
zero.  Thus its long exact sequence contains
\begin{equation}
  HC^{j-2}_m\longrightarrow HC^j_m\longrightarrow HH^j_m.
  \label{eq:polygon-cyclic-sequence}
\end{equation}
For $m>2$, the last space vanishes at every $j$ by
\eqref{eq:ordinary-bar-bidegrees}.  Starting with zero in negative
cochain degrees and inducting in $j$ proves $HC^j_m=0$ for every
$j$ and $m>2$.  The single-trace summand in Chevalley--Eilenberg
arity $k$ is $HC^{k-1}_m$, not $HC^k_m$.
Deleting the one-letter trace in the traceless complex cannot change
cohomology in arities three and four: their incoming and outgoing
maps involve only arities two through five.  Hence the one-trace
summands of both groups in \eqref{eq:all-order-cohomology} vanish.

\emph{Two traces.}
At arity four the only remaining word type is $(01)(23)$.
Write its coefficient as $F(r,q,P)$ on the momenta
$(r,P-r,q,-P-q)$.  Alternation gives
\begin{align}
 F(P-r,q,P)&=-F(r,q,P),\notag\\
 F(r,-P-q,P)&=-F(r,q,P),\notag\\
 F(q,r,-P)&=F(r,q,P).
 \label{eq:pair-exchanges}
\end{align}
The coefficient of an oriented three-letter trace times the other
two-letter trace in $dF$ says
\begin{equation}
 F(r,q,P)+F(s,q,P)+F(P-r-s,q,P)=0.
\end{equation}
Using \eqref{eq:pair-exchanges}, this is precisely
\begin{equation}
 F(r+s,q,P)=F(r,q,P)+F(s,q,P).
 \label{eq:pair-additivity}
\end{equation}
Polynomial additivity implies linearity in $r$.  The first exchange
then forces that linear functional to annihilate $P$.  The two
coordinates of $P$ are relatively prime in their polynomial ring;
therefore a polynomial linear functional annihilating $P$ is a
polynomial multiple of $[rP]$.  Applying the same argument to $q$
gives
\begin{equation}
 F(r,q,P)=[rP][qP]h(P),\qquad h(-P)=h(P).
\end{equation}
Conversely this expression satisfies every coefficient of $dF=0$.
At total momentum degree $2w$, the polynomial $h$ has degree
$2w-4$, so its evenness is automatic.  There is no arity-three
two-trace word on traceless matrices.  These cocycles are therefore
not boundaries.  Their $SL_2$ representation and dimension are those
of binary forms of degree $2w-4$.  This proves all assertions.
\end{proof}

\subsection{Induction and finite-rank specialization}

\begin{theorem}[All-order traceless quantization]
\label{thm:all-orders-traceless}
Take the defining trace form on $\fsl_N$ and the flat heat-kernel
scheme.  For each of the two complex solutions of
\begin{equation}
 a^2=\epsilon U\frac{N^2}{256\pi^4},
 \label{eq:all-order-leading-root}
\end{equation}
the first-order cubic $J_1=aF_D$ has a formal local quantization
\begin{equation}
 S_{\mathrm{cl}}+\hbar J_1+\sum_{w\ge2}\hbar^w I_w,
 \qquad I_w\in\Vht,
 \label{eq:all-order-action}
\end{equation}
with the prescribed $U(2)$ characters.  In the formal trace algebra
each $I_w$ is the unique invariant cubic solving the order-$w$
equation.  The complete identities specialize to every integer
$N\ge2$.  At $N=2$, $J_1$ is identically zero; the conclusion there
is existence and does not infer intrinsic uniqueness from the formal
trace calculation.
\end{theorem}

\begin{proof}
The two-loop calculation supplies $J_1$ and $I_2$ satisfying the local
equation through order two.  Suppose the local equations below order
$w\ge3$ have been solved.  The leading remaining obstruction is
classically closed.  One way to see this is to use the finite-scale
quantum Batalin--Vilkovisky identity.  With $Q$ the free differential,
$\Delta_L$ the heat-regularized contraction, and $S[L]$ the effective
interaction, its curvature satisfies
\begin{gather}
 O[L]=QS[L]+\hbar\Delta_LS[L]
                  +\tfrac12\{S[L],S[L]\}_L,\notag\\
 \bigl(Q+\hbar\Delta_L+\{S[L],-\}_L\bigr)O[L]=0.
 \label{eq:finite-scale-bianchi}
\end{gather}
At the first unsolved order the lower curvatures vanish.  Passing to
the local contact, or equivalently applying the local obstruction
construction of~\cite{CostelloBV}, gives classical closedness.
The compactified boundary factorization of~\cite[Proposition
3.22]{WW} identifies this local equation with the effective master
equation at positive scale.

Lemma~\ref{lem:full-wedge-contacts} puts the actual local obstruction
in $\Vht$, so its Taylor cochain is an element of
$C^{0,4}_{\fsl_N}$.  The determinant character gives total momentum
degree $2w$.  Its $SU(2)$ invariance is $SL_2$ invariance because the
momentum coefficients are holomorphic polynomials.  By
\eqref{eq:all-order-invariant-vanishing}, its invariant cohomology
class is zero.  Complete reducibility in each finite degree gives
an invariant primitive.  The chain identification with $\Vht$
turns this primitive into the required local counterterm $I_w$.

To verify uniqueness within the stated formal sector, three
conserved momenta leave two independent momentum vectors.
Their invariant polynomial ring is generated by their determinant.
The cubic Taylor cochain at order $w$ is therefore a scalar multiple
of
\begin{equation}
 [p_0p_1]^w
 \begin{cases}
  \tr\bigl(X_0[X_1,X_2]\bigr),&w\text{ even},\\
  \tr\bigl(X_0\{X_1,X_2\}\bigr),&w\text{ odd}.
 \end{cases}
 \label{eq:all-order-cubic-shape}
\end{equation}
Here $\{X_1,X_2\}=X_1X_2+X_2X_1$ is the matrix anticommutator.
The parity choices make the entire cochain alternating.  There is
no invariant two-cochain of positive momentum degree, since two
conserved momenta are $p,-p$ and a single vector has no
nonconstant $SL_2$-invariant polynomial.  Together with
$H^3_{\deg=2w}=0$, this proves that the differential is injective
on the invariant cubic space for $w\ge2$.

For specialization, choose $a/N$ to be either fixed square root in
\eqref{eq:all-order-leading-root}.  The inverse trace pairing on
traceless matrices has the exact completeness identity
\begin{equation}
 \sum_b\tr(XT_b)\tr(T^bY)
   =\tr(XY)-\frac1N\tr(X)\tr(Y),
 \label{eq:all-order-completeness}
\end{equation}
where $T_b,T^b$ are trace-dual bases of $\fsl_N$.
Each graph consequently has coefficients in
$\mathbb C[N,N^{-1}]$, its graph periods being scalar constants.
The formal trace differential has rational coefficients independent
of $N$.  At each fixed momentum degree, solving on its image
therefore introduces only nonzero rational denominators, not new
polynomials in $N$ in denominators.  Induction constructs the
entire local identity over $\mathbb C[N,N^{-1}]$.
Evaluation on traceless matrices at a positive integer $N$ commutes
with the classical differential, tensor contractions, and graph
composition.  Hence every specialized identity holds, whether or
not additional trace relations exist at that rank.

For $N=2$, the identity
$X_1X_2+X_2X_1=\tr(X_1X_2)\mathbf1_2$ on traceless matrices
annihilates every odd-order cubic in
\eqref{eq:all-order-cubic-shape}.  In particular $J_1=0$.
The even-order specialized counterterms still satisfy the full
local equations.  No inference about the intrinsic cohomology
of the small-rank complex is used.
\end{proof}

\begin{remark}[What is and is not classified]
The uniqueness statement fixes the heat-kernel scheme, the leading
root, and the full-wedge equivariant cubic representatives.  It is
not a statement that arbitrary local field redefinitions or arbitrary
renormalization schemes are identical.  The formal series need not
converge analytically in $\hbar$.  Its coefficients are local and of
finite derivative order at each order, which is the perturbative
meaning of quantization used here.
\end{remark}

\subsection{A finite evaluation determining every higher cubic}

The all-order theorem also gives an exact extraction formula for the
counterterm coefficient.  This is useful independently of any
choice of basis for four-cochains.  For $N\ge4$, evaluate an
order-$w$ equation on the color matrices
\begin{equation}
  (X_0,X_1,X_2,X_3)=(E_{12},E_{23},E_{34},E_{41})
\end{equation}
and on the conserved momenta
\begin{equation}
  (p_0,p_1,p_2,p_3)=((1,0),(0,1),(1,1),(-2,-2)).
\end{equation}
Here $E_{ab}$ is the matrix with its only nonzero entry equal to one
in row $a$ and column $b$.  The differential of the unscaled cubic
in \eqref{eq:all-order-cubic-shape} evaluates to
\begin{align}
 -[p_2p_3]^w+[p_1p_2]^w+[p_3p_0]^w-[p_0p_1]^w
   &=2^w+(-1)^w-1\notag\\
   &=\begin{cases}2^w,&w\text{ even},\\
                    2^w-2,&w\text{ odd}.
      \end{cases}
 \label{eq:all-order-extraction-denominator}
\end{align}
It is nonzero for every $w\ge2$.  Dividing the evaluated right-hand
side of the cochain equation $d i_w=\text{known lower-order terms}$
by this integer determines its unique invariant cubic primitive.
This extraction does not prove closedness; the master-equation
identity and the cohomology calculation supply that logically prior
step.  Explicit four-loop periods and the existing independent
matrix calculations are recorded in~\cite{RepoHigher}.
